\section*{Limitations}

\textbf{Episode scope.} Test episodes are generated by GPT-4.1 and carry one future question each. Evaluating multiple future questions per episode will test how question diversity affects the Statedness Bias.

\textbf{Statedness.} Stem matching approximates whether a concept is stated. Semantic annotation should test this definition on paraphrases and indirect mentions.

\textbf{Answer grading.} Answer evaluation relies on a blind semantic judge, with 85.1\% agreement with string matching. Human grading of the disagreements is needed to assess their effects on comparisons between admission policies.

\textbf{Candidate supply.} Candidates are supplied upstream, so the evaluation isolates admission from extraction, update and retrieval. The next step is to evaluate how these operations interact with the Statedness Bias in full memory systems.

\textbf{Serving.} Serving measurements cover Qwen3-8B in one configuration. Broader serving tests should measure throughput, answer accuracy, and agreement with offline admission across models and configurations.

\section*{Ethical Considerations}

\textbf{Artifacts.} We use the six open-weight models under their licenses and release the episodes, candidate pools and code under CC BY 4.0 for research use. GPT-4.1 generated the released episodes, and an LLM generated the unstated distractors. The episodes name invented characters by first name and contain no personal data.

\textbf{Use of AI assistants.} AI assistants were used to draft and edit text and to write code; the authors verified all outputs and numbers.
